\section{Conclusion}

In this paper, we delineate \app{} into four distinct categories based on the primary focus of queries, each facing unique challenges and thus requiring tailored solutions, as illustrated in Figure~\ref{fig:conclusion}. For queries related to static common knowledge, deploying a general LLM through a Chain of Thought methodology is effective. For \firstlevel, the main challenge involves pinpointing the exact location of facts within a database, thus making a basic RAG the method of choice. In the case of \secondlevel, which require the collation of multiple related facts, iterative RAG and RAG implementations on graph or tree structures are preferred for their ability to concurrently retrieve individual facts and interconnect multiple data points. When extensive data linkage is necessary, text-to-SQL techniques prove indispensable, leveraging database tools to facilitate external data searches. For \thirdlevel, advancements through prompt tuning and CoT prompting are critical to enhance LLMs' compliance with external directives. The most formidable are \forthlevel, which demand the autonomous synthesis of problem-solving approaches from extensive data sets. Here, offline learning, in-context learning , and fine-tuning emerge as vital methodologies.

\begin{figure}
    \centering
    \includegraphics[width=0.4\textwidth]{contents/images/three_type_input.png}
    \caption{Three ways to inject specific domain data into an LLM: a) extracting part of the domain data based on the query as context input for the LLM, b) training a smaller model with specific domain data, which then guides the integration of external information subsequently input into the LLM, and c) directly using external domain knowledge to fine-tune a general large language model to become a domain-expert model.}
    \label{fig:three_type_input}
\end{figure}

Prior to the development of a targeted LLM application, as domain experts, we must acquire an in-depth understanding of the intended task, ascertain the complexity level of the associated queries, and select corresponding technological approaches for resolution. These methods principally inject knowledge into LLMs via three mechanisms, as depicted in Figure~\ref{fig:three_type_input}: a) extracting part of the domain data based on the query as context input for the LLM, b) training a smaller model with specific domain data, which then guides the integration of external information subsequently input into the LLM, and c) directly using external domain knowledge to fine-tune a general large language model to become a domain-expert model. These strategies differ in their requirements for data volume, training duration, and computational resources, escalating respectively. Knowledge injection through context provides better interpretability and stability but faces limitations due to the finite context window and potential information loss in the middle~\cite{liu2024lost}, ideally suited for scenarios where data can be succinctly explained in shorter texts. However, this method challenges the model's retrieval capabilities and knowledge extraction ability. The small model approach offers the advantage of reduced training times and the capacity to assimilate considerable amounts of data, yet its efficacy is contingent upon the model's capabilities, potentially capping the LLM’s performance for more complex tasks and incurring additional training costs with data increasing. Fine-tuning facilitates the utilization of large model capacities with extensive domain-specific data, yet its impact on the LLM strongly depends on the design of the data used. Employing out-of-domain factual data for fine-tuning may inadvertently lead to the generation of more erroneous outputs by the LLM, while also risking the loss of previously known domain knowledge and the neglect of unencountered tasks during fine-tuning~\cite{fine_tune_new_24, sun2024amuro}. Therefore, choosing an appropriate data injection strategy into the LLM requires a thorough understanding of one's data sources and judicious decision-making based on this insight.

Moreover, in practical scenarios, \app typically involves a combination of diverse query types, necessitating developers to engineer a routing pipeline that integrates multiple methodologies to effectively tackle these multifaceted challenges.